% ============================
% checkoff.tex - TA check-offs (loaded by header.tex)
% ----------------------------
% \begin{checkoff}[Short title]
%   What to show the TA.
% \end{checkoff}
%     A numbered call-out wherever work must be demonstrated to a TA or the
%     instructor. It works in running text and inside lists, and may contain
%     inline code (it is an environment, not a command, so \mintinline and
%     \lstinline work inside). The box carries a checkbox the reader can tick:
%       PDF (normal and large print)  a tagged form field (l3pdffield); it
%                                     prints as an empty box
%       HTML                          <input type="checkbox">, remembered in the
%                                     browser (script in se423.cfg)
%       EPUB                          <input type="checkbox"> (no script)
%     The short title (plain text) is what the summary lists; without it the
%     summary says "Check-Off n".
%
% \CheckOffSummary
%     The list of every check-off of the document (put it at the end). Each
%     line has the same checkbox as its call-out (ticking one ticks the other,
%     except in EPUB, which lists links only) and links back to the call-out.
%
% \CheckOffCount, \CheckOffNotice
%     The number of check-offs in the document, and a ready-made line with it
%     for the title box (end \AssignmentText with "\par \CheckOffNotice").
%
% The list and the count come from the .aux file, so they are right from the
% second LaTeX run on (latexmk, make4ht and tex4ebook all run more than once).
% Colours: only `highlight` (header.tex) and greys, and never colour alone: every
% call-out starts with the words "TA Check-Off".
% ============================
\newcounter{SEcheckoff}
\makeatletter
% --- list of check-offs, collected through the .aux file --------------------
\gdef\SE@checkofflist{}
\gdef\SE@checkoffcount{0}
% Written to the .aux by \CheckOff: {number}{short title}
\newcommand{\SE@checkoffentry}[2]{%
  \g@addto@macro\SE@checkofflist{\SE@checkoffline{#1}{#2}}%
  \gdef\SE@checkoffcount{#1}%
}
\newcommand{\CheckOffCount}{\SE@checkoffcount}
% A line for the title box (\AssignmentText): how many check-offs to expect.
\newcommand{\CheckOffNotice}{%
  \textcolor{highlight}{TA Check-Offs in this document: \CheckOffCount}
  (marked ``TA Check-Off'' in the text and listed at the end)%
}
\newcommand{\SE@checkoffrecord}[1]{%
  \immediate\write\@auxout{\string\SE@checkoffentry{\theSEcheckoff}{\unexpanded{#1}}}%
}
\AtEndDocument{%
  \ifnum\SE@checkoffcount=\value{SEcheckoff}\else
    \@latex@warning@no@line{Check-off list changed. Rerun to get the summary right}%
  \fi
}
\makeatother
%
\ifSEhtml
  % ----- HTML / EPUB ---------------------------------------------------------
  \makeatletter
  \NewDocumentEnvironment{checkoff}{O{}}{%
    \refstepcounter{SEcheckoff}%
    \SE@checkoffrecord{#1}%
    \ifvmode\IgnorePar\fi\EndP
    % Link target made by tex4ht, so the summary's links also resolve when the
    % document is split over several files (EPUB).
    \Link{}{checkoff-\theSEcheckoff}\EndLink
    \HCode{<div class="checkoff" role="note"
      aria-labelledby="checkoff-label-\theSEcheckoff"><p class="checkoff-head"><input
      type="checkbox" id="checkoff-box-\theSEcheckoff" data-checkoff="\theSEcheckoff" />
      <label id="checkoff-label-\theSEcheckoff" for="checkoff-box-\theSEcheckoff">TA
      Check-Off \theSEcheckoff}%
    \if\relax\detokenize{#1}\relax\else\HCode{: }#1\fi
    \HCode{</label></p><div class="checkoff-body">}%
    \par\ShowPar
  }{%
    \ifvmode\IgnorePar\fi\EndP
    \HCode{</div></div>}%
    \par
  }
  \newcommand{\SE@checkoffline}[2]{%
    \HCode{<li>}%
    \ifdefined\SEebook\else
      \HCode{<input type="checkbox" data-checkoff="#1"
        aria-labelledby="checkoff-sum-#1" /> }%
    \fi
    \Link{checkoff-#1}{checkoff-sum-#1}Check-Off #1%
    \if\relax\detokenize{#2}\relax\else: #2\fi
    \EndLink
    \HCode{</li>}%
  }
  \newcommand{\CheckOffSummary}{%
    \ifvmode\IgnorePar\fi\EndP
    \HCode{<ul class="checkoff-summary">}%
    \SE@checkofflist
    \HCode{</ul>}%
    \par
  }
  \makeatother
\else
  % ----- PDF -----------------------------------------------------------------
  \usepackage{l3pdffield}
  \ExplSyntaxOn
  % An empty box, and the same box with a check mark drawn in it (no font
  % needed, so the mark looks the same in every viewer and in print).
  \pdffield_appearance:nn { SEcheckoff/Off }
    {
      \begin{tikzpicture}[x=1em,y=1em]
        \draw[line~width=0.9pt,rounded~corners=0.12em,fill=white] (0,0) rectangle (1,1);
      \end{tikzpicture}
    }
  \pdffield_appearance:nn { SEcheckoff/Yes }
    {
      \begin{tikzpicture}[x=1em,y=1em]
        \draw[line~width=0.9pt,rounded~corners=0.12em,fill=white] (0,0) rectangle (1,1);
        \draw[line~width=1.6pt,line~cap=round,line~join=round]
          (0.2,0.52) -- (0.42,0.24) -- (0.82,0.8);
      \end{tikzpicture}
    }
  % #1 number, #2 accessible name (both fully expanded: they name the field)
  \NewDocumentCommand{\SEcheckoffbox}{m m}
    { \use:e { \exp_not:N \SE_checkoff_box:nn {#1} {#2} } }
  \cs_new_protected:Npn \SE_checkoff_box:nn #1#2
    {
      \leavevmode
      \raisebox{-0.18em}
        {
          % PDF/UA-2 8.10.2.3: a widget without a visible label of its own
          % needs a Contents entry (the description assistive technology reads).
          \pdf_string_from_unicode:nnN { utf16/hex } {#2} \l_tmpa_str
          \pdfannot_dict_put:nne { widget } { Contents } { \l_tmpa_str }
          \pdffield_checkbox:n
            {
              name = checkoff#1,
              altname = {#2},
              appearance = SEcheckoff,
              setannotflags = Print, % an unticked box still prints
              width = 1em, height = 1em, depth = 0pt,
            }
        }
    }
  \ExplSyntaxOff
  \makeatletter
  \NewDocumentEnvironment{checkoff}{O{}}{%
    \par
    \refstepcounter{SEcheckoff}%
    \SE@checkoffrecord{#1}%
    \begin{tcolorbox}[
        enhanced, breakable=false,
        colback=highlight!6!white,
        colframe=highlight,
        boxrule=0.6pt, leftrule=3mm,
        arc=1mm,
        left=2.5mm, right=2.5mm, top=1.2mm, bottom=1.2mm,
        before skip=2mm, after skip=2mm,
        width=\linewidth,
    ]
    \raggedright
    \hypertarget{checkoff.\theSEcheckoff}{}%
    \SEcheckoffbox{\theSEcheckoff}{TA check-off \theSEcheckoff\space done}\enspace
    {\bfseries\color{highlight}TA Check-Off \theSEcheckoff
      \if\relax\detokenize{#1}\relax\else: #1\fi}\par
    \vspace{0.6mm}%
    \ignorespaces
  }{%
    \par
    \end{tcolorbox}
    \par
  }
  \newcommand{\SE@checkoffline}[2]{%
    \item \SEcheckoffbox{#1}{TA check-off #1 done}\enspace
      \hyperlink{checkoff.#1}{Check-Off #1\if\relax\detokenize{#2}\relax\else: #2\fi}%
  }
  \newcommand{\CheckOffSummary}{%
    \par
    \ifnum\SE@checkoffcount>0
      \begin{itemize}
        \renewcommand{\labelitemi}{}%
        \SE@checkofflist
      \end{itemize}
    \fi
  }
  \makeatother
\fi
